\documentclass[../../main.tex]{subfiles}
\begin{document}

\section{环面和小丑帽的基本群}\label{section:ch11-环面和小丑帽的基本群}

在本节中，我们将应用上一节的结果来计算两个空间的基本群，其一是已知的，而另一却未知。所用到的技巧对于后续讨论具有重要意义。

\begin{theorem}\label{theorem:ch11-torus-presentation}
环面的基本群有一个由两个生成元 $\alpha$，$\beta$ 及单一关系 $\alpha\beta\alpha^{-1}\beta^{-1}$ 组成的表示。
\end{theorem}

\par
\begin{figure}[htbp]
\centering
\includegraphics[width=0.6\textwidth]{image/fig-73-1.png}
\caption{}\label{figure:ch11-fig-73-1}
\end{figure}

\begin{proof}
设 $X=S^1\times S^1$ 为环面，$h:I^2\to X$ 为标准覆叠映射 $p\times p:\mathbb{R}\times\mathbb{R}\to S^1\times S^1$ 的限制。设 $p$ 为 $\operatorname{Bd}I^2$ 中的点 $(0,0)$，$a=h(p)$，$A=h(\operatorname{Bd}I^2)$。则定理 72.1 的假设条件成立。

空间 $A$ 为两个圆周的束，所以 $A$ 的基本群为自由群。事实上，若在 $\operatorname{Bd}I^2$ 中令 $a_0$ 为道路 $a_0(t)=(t,0)$，$b_0$ 为道路 $b_0(t)=(0,t)$，则道路 $\alpha=h\circ a_0$ 和道路 $\beta=h\circ b_0$ 都是 $A$ 中的回路，并且 $[\alpha]$ 和 $[\beta]$ 构成 $\pi_1(A,a)$ 的一个自由生成元组。见图 73.1。


设 $a_1$ 和 $b_1$ 为 $\operatorname{Bd}I^2$ 中的道路 $a_1(t)=(t,1)$ 和 $b_1(t)=(1,t)$。考虑由
\begin{align*}
f=a_0*(b_1*(\bar{a}_1*\bar{b}_0))
\end{align*}
定义的 $\operatorname{Bd}I^2$ 中的回路 $f$，则 $f$ 表示 $\pi_1(\operatorname{Bd}I^2,p)$ 的一个生成元，并且回路 $g=h\circ f$ 等于乘积 $\alpha*(\beta*(\bar{\alpha}*\bar{\beta}))$。根据定理 72.1，$\pi_1(X,a)$ 是由自由生成元 $[\alpha]$ 和 $[\beta]$ 生成的自由群对于含有元素 $[\alpha][\beta][\alpha]^{-1}[\beta]^{-1}$ 的最小正规子群的商群。
\end{proof}

\begin{corollary}\label{corollary:ch11-torus-abelian}
环面的基本群是秩为 2 的自由阿贝尔群。
\end{corollary}

\begin{proof}
设 $G$ 是以 $\alpha$，$\beta$ 为生成元的自由群，$N$ 是含有元素 $\alpha\beta\alpha^{-1}\beta^{-1}$ 的最小正规子群。因为这个元素为一个交换子，所以 $N$ 包含于 $G$ 的交换子子群 $[G,G]$ 中。另一方面，因为 $G/N$ 是由所有陪集 $\alpha N$ 及 $\beta N$ 所生成的，所以它的元素可交换，即 $G/N$ 是阿贝尔群。因此 $N$ 包含着 $G$ 的交换子子群。

根据定理 69.4，$G/N$ 是一个秩为 2 的自由阿贝尔群。
\end{proof}

\begin{definition}[n-叠小丑帽]\label{definition:ch11-dunce-cap}
设 $n>1$ 为正整数。设 $r:S^1\to S^1$ 是以 $2\pi/n$ 为旋转角的旋转，它将点 $(\cos\theta,\sin\theta)$ 映为点 $(\cos(\theta+2\pi/n),\sin(\theta+2\pi/n))$。在单位球 $B^2$ 中将 $S^1$ 中的每一个点 $x$ 与点 $r(x)$，$r^2(x)$，$\cdots$，$r^{n-1}(x)$ 等同起来，得到的商空间记为 $X$。我们将要证明 $X$ 是一个 Hausdorff 空间，并且将这个空间称为 $n$-叠小丑帽（$n$-fold dunce cap）。
\end{definition}

设 $\pi:B^2\to X$ 为前面定义的商映射。我们证明 $\pi$ 是闭映射。为此，只需证明如果 $C$ 是 $B^2$ 中的闭集，则 $\pi^{-1}\pi(C)$ 为 $B^2$ 中的闭集，然后由商映射的定义可见 $\pi(C)$ 为 $X$ 中的闭集。设 $C_0=C\cap S^1$，它是 $B^2$ 中的闭集。集合 $\pi^{-1}\pi(C)$ 等于 $C$ 与 $r(C_0)$，$r^2(C_0)$，$\cdots$，$r^{n-1}(C_0)$ 的并，由于 $r$ 为同胚，上述每一个集合都是 $B^2$ 中的闭集。因此，$\pi^{-1}\pi(C)$ 为 $B^2$ 中的闭集。

由于 $\pi$ 连续，所以 $X$ 是紧致的。曾经在第 31 节中作为一个习题给出的下述引理说明 $X$ 是一个 Hausdorff 空间。

\begin{lemma}\label{lemma:ch11-closed-quotient-regular}
设 $\pi:E\to X$ 是一个闭的商映射。若 $E$ 是正规的，则 $X$ 也是正规的。
\end{lemma}

\par
\begin{figure}[htbp]
\centering
\includegraphics[width=0.5\textwidth]{image/fig-73-2.png}
\caption{}\label{figure:ch11-fig-73-2}
\end{figure}

\begin{proof}
假设 $E$ 为正规空间。由于 $E$ 中的单点集为闭集，所以 $X$ 中的单点集为闭集。设 $A$ 和 $B$ 为 $X$ 中的无交闭集，则 $\pi^{-1}(A)$ 和 $\pi^{-1}(B)$ 为 $E$ 中的无交闭集。选取 $U$ 和 $V$ 为 $E$ 中分别包含 $\pi^{-1}(A)$ 和 $\pi^{-1}(B)$ 的无交开集。当然，$\pi(U)$ 和 $\pi(V)$ 为分别包含 $A$ 和 $B$ 的无交开集这样一个猜测对我们很有诱惑力，然而，这个猜测是不成立的。因为它们未必是开集（$\pi$ 未必是开映射），也未必是无交的。见图 73.2。


正因为如此，我们才按以下方式给出证明。设 $C=E-U$ 和 $D=E-V$。由于 $C$ 和 $D$ 是 $E$ 中的两个闭集，所以 $\pi(C)$ 和 $\pi(D)$ 是 $X$ 中的两个闭集。由于 $C$ 不包含 $\pi^{-1}(A)$ 中的点，所以 $\pi(C)$ 与 $A$ 无交。因此 $U_0=X-\pi(C)$ 为 $X$ 中包含 $A$ 的开集。类似地，$V_0=X-\pi(D)$ 为 $X$ 中包含 $B$ 的开集。进而，$U_0$ 与 $V_0$ 无交。这是因为，如果 $x\in U_0$，则 $\pi^{-1}(x)$ 与 $C$ 无交，因此它包含于 $U$ 中。类似地，如果 $x\in V_0$，则 $\pi^{-1}(x)$ 包含于 $V$ 中。由于 $U$ 和 $V$ 无交，所以 $U_0$ 与 $V_0$ 也无交。
\end{proof}

事实上，2-叠小丑帽是我们曾见过的一个空间，它同胚于射影平面 $P^2$。为了验证这个事实，我们回顾一下 $P^2$ 的定义。它是将 $S^2$ 中每一对对径点 $x$ 和 $-x$ 等同所得到的商空间。设 $p:S^2\to P^2$ 为相应的商映射。用 $i$ 表示 $B^2$ 与 $S^2$ 中上半球面之间的标准同胚，即
\begin{align*}
i(x,y)=(x,y,(1-x^2-y^2)^{1/2}),
\end{align*}
再将其与 $p$ 复合。我们得到一个映射 $\pi:B^2\to P^2$，这个映射是连续的、闭的、满的。在 $\operatorname{Int}B^2$ 上它是一个单射，并且对任何一点 $x\in S^1$，它将 $x$ 和 $-x$ 映为同一点。因此它诱导出一个 2-叠小丑帽到 $P^2$ 的同胚。

从 $P^2$ 的基本群的计算中你便能猜到关于 $n$-叠小丑帽的基本群的结论。

\begin{theorem}\label{theorem:ch11-dunce-cap-fundamental-group}
$n$-叠小丑帽的基本群是一个 $n$ 阶循环群。
\end{theorem}

\par
\begin{figure}[htbp]
\centering
\includegraphics[width=0.5\textwidth]{image/fig-73-3.png}
\caption{}\label{figure:ch11-fig-73-3}
\end{figure}

\begin{proof}
设 $h:B^2\to X$ 为商映射，其中 $X$ 为 $n$-叠小丑帽。令 $A=h(S^1)$，$p=(1,0)\in S^1$，$a=h(p)$。则 $h$ 将 $S^1$ 中的从 $p$ 到 $r(p)$ 的弧 $C$ 映到 $A$ 上，它将 $C$ 的两个端点映成一个点，但在其他地方为单射。因此，$A$ 同胚于圆周，于是它的基本群是无限循环群。事实上，如果 $\gamma$ 是 $S^1$ 中从 $p$ 到 $r(p)$ 由
\begin{align*}
\gamma(t)=(\cos(2\pi t/n),\sin(2\pi t/n))
\end{align*}
定义的道路，则 $\alpha=h\circ\gamma$ 表示 $\pi_1(A,a)$ 的一个生成元。参见图 73.3。


于是回路
\begin{align*}
f=\gamma*((r\circ\gamma)*((r^2\circ\gamma)*\cdots*(r^{n-1}\circ\gamma)))
\end{align*}
的同伦类生成 $\pi_1(S^1,p)$。由于 $h(r^m(x))=h(x)$ 对于所有的 $x$ 及 $m$ 成立，回路 $h\circ f$ 等于 $n$-重乘积 $\alpha*(\alpha*(\cdots*\alpha))$。因此定理成立。
\end{proof}

\subsection*{习题}

\begin{exercise}\label{exercise:ch11-73-01}
试给出其基本群同胚于下述群的空间。（其中 $\mathbb{Z}/n$ 表示模 $n$ 整数加群。）
\begin{enumerate}[(a)]
\item $\mathbb{Z}/n\times\mathbb{Z}/m$。
\item $\mathbb{Z}/n_1\times\mathbb{Z}/n_2\times\cdots\times\mathbb{Z}/n_k$。
\item $\mathbb{Z}/n*\mathbb{Z}/m$。（参见第 71 节习题 2。）
\item $\mathbb{Z}/n_1*\mathbb{Z}/n_2*\cdots*\mathbb{Z}/n_k$。
\end{enumerate}
\end{exercise}

\begin{exercise}\label{exercise:ch11-73-02}
证明下述定理：

定理 设 $G$ 是一个有限表示的群，则存在一个紧致的 Hausdorff 空间 $X$，其基本群同胚于 $G$。

证明：假设 $G$ 有一个由 $n$ 个生成元和 $m$ 个关系组成的表示。设 $A$ 为 $n$ 个圆周的束，由 $A$ 与 $m$ 个单位球的同胚像 $B_1,\cdots,B_m$ 的并以及借助于连续映射 $f:\bigcup\operatorname{Bd}B_i\to A$ 所得到的贴附空间记作 $X$。
\begin{enumerate}[(a)]
\item 证明 $X$ 为 Hausdorff 空间。
\item 对于 $m=1$ 的情形证明定理。
\item 应用下一个习题中表述的代数学中的结论，对 $m$ 进行归纳。
\end{enumerate}

在这个习题中提到的构造空间 $X$ 的方法是代数拓扑中的一个标准方法。空间 $X$ 叫做 2 维 CW 复形（CW complex）。
\end{exercise}

\begin{exercise}\label{exercise:ch11-73-03}
引理 设 $f:G\to H$ 和 $g:H\to K$ 都是同态，$f$ 是一个满射。若 $x_0\in G$ 并且 $\ker g$ 为 $H$ 中包含 $f(x_0)$ 的最小正规子群，则 $\ker(g\circ f)$ 是 $G$ 中含有 $\ker f$ 和 $x_0$ 的最小正规子群 $N$。

证明 $f(N)$ 正规，并且据此证明 $\ker(g\circ f)=f^{-1}(\ker g)\subset f^{-1}f(N)=N$。
\end{exercise}

\begin{exercise}\label{exercise:ch11-73-04}
证明：习题 2 中所构造的空间实际上是可度量化的。[提示：商映射是一个完全映射。]
\end{exercise}

\end{document}